\section{Experiments and Observations}
\label{sec:results}

All numbers are means over the three test seeds, with the per-seed range in brackets. They are
descriptive; with three seeds we report neither significance tests nor equivalence.

\subsection{A loss of learning ability on a comparable new task}
\label{sec:phenomenon}

\begin{table}[t]
\caption{Probe AUC (mean held-out probe accuracy over 200 steps) from three starting points, base Adam.
Differences are late minus early, per seed paired; brackets give the seed range.}
\label{tab:phenomenon}
\centering\small
\begin{tabular}{lccccc}
\toprule
Condition & fresh & early & late & late$-$early (AUC) & late$-$early (final acc.) \\
\midrule
random-teacher ($C{=}4$) & 0.788 & 0.727 & 0.483 & -0.245 [-0.315, -0.198] & -0.177 [-0.258, -0.133] \\
label-permutation ($C{=}10$) & 0.568 & 0.536 & 0.332 & -0.204 [-0.282, -0.151] & -0.202 [-0.246, -0.137] \\
\bottomrule
\end{tabular}

\end{table}

\begin{figure}[t]
\centering
\begin{tikzpicture}
\begin{axis}[width=0.43\linewidth, height=4.2cm, title={random-teacher ($C{=}4$)}, xlabel={probe step},
  ylabel={held-out probe acc.}, ymin=0, ymax=1, legend pos=south east, legend style={font=\scriptsize},
  tick label style={font=\scriptsize}, label style={font=\small}, title style={font=\small}]
\addplot[thick, black] table[x=k, y=mean] {figures/curve_RT_fresh.dat}; \addlegendentry{fresh}
\addplot[thick, blue, dashed] table[x=k, y=mean] {figures/curve_RT_early.dat}; \addlegendentry{early}
\addplot[thick, red, dotted] table[x=k, y=mean] {figures/curve_RT_late.dat}; \addlegendentry{late}
\end{axis}
\end{tikzpicture}\hfill
\begin{tikzpicture}
\begin{axis}[width=0.43\linewidth, height=4.2cm, title={label-permutation ($C{=}10$)}, xlabel={probe step},
  ymin=0, ymax=1, tick label style={font=\scriptsize}, label style={font=\small}, title style={font=\small}]
\addplot[thick, black] table[x=k, y=mean] {figures/curve_LP_fresh.dat};
\addplot[thick, blue, dashed] table[x=k, y=mean] {figures/curve_LP_early.dat};
\addplot[thick, red, dotted] table[x=k, y=mean] {figures/curve_LP_late.dat};
\end{axis}
\end{tikzpicture}
\caption{Learning curves on the same never-trained probe task from the fresh, early (after task 0) and
late (after task 49) starting points, with identical batches, dropout masks and update budget (mean of
three seeds; per-seed curves are in \texttt{paper/figures/*.dat}).}
\label{fig:curves}
\end{figure}

Table~\ref{tab:phenomenon} and Figure~\ref{fig:curves} show that the late network learns the probe task
more slowly and to a lower final accuracy than the early network, in both conditions and on every seed.
The late-minus-early AUC is \RTLateMinusEarlyAuc{} [\RTLateMinusEarlyAucMin, \RTLateMinusEarlyAucMax]
for random-teacher and \LPLateMinusEarlyAuc{} [\LPLateMinusEarlyAucMin, \LPLateMinusEarlyAucMax] for
label-permutation. The first-50-step accuracy differs by \RTLateMinusEarlyEarlyAcc{} and
\LPLateMinusEarlyEarlyAcc{}. By the pre-declared rule the phenomenon is established in both conditions.
The early network is itself slightly worse than fresh (\RTEarlyMinusFreshAuc{} and
\LPEarlyMinusFreshAuc{}). The fresh comparison reflects the effect of past training on the weights
together with a fresh optimizer. It is a diagnostic reference, not a state-only comparison. The gap also
appears with the separately tuned $\beta_1{=}\beta_2{=}0.9$ Adam (\RTAdamEqBetaLateMinusEarlyAuc{},
\LPAdamEqBetaLateMinusEarlyAuc{}) and SGD with momentum (\RTSgdLateMinusEarlyAuc{},
\LPSgdLateMinusEarlyAuc{}). In this setting it is therefore not specific to Adam's default $\beta$ values.

\subsection{State-only interventions from the late checkpoint}
\label{sec:interventions}

\begin{table}[t]
\caption{Interventions on the late checkpoint. $\Delta$: probe-AUC difference to \keep{} and to the
intervention's update-norm control, mean [seed range]. Update norms and descriptive statistics are in
Appendix~\ref{app:descriptive}. $^\dagger$Mis-corrected artifact controls (counter not reset), not
candidate methods.}
\label{tab:interventions}
\centering\scriptsize\setlength{\tabcolsep}{3pt}
\begin{tabular}{l cc cc}
\toprule
 & \multicolumn{2}{c}{random-teacher ($C{=}4$)} & \multicolumn{2}{c}{label-permutation ($C{=}10$)} \\
\cmidrule(lr){2-3}\cmidrule(lr){4-5}
Arm & $\Delta$ vs.\ keep & $\Delta$ vs.\ norm ctrl. & $\Delta$ vs.\ keep & $\Delta$ vs.\ norm ctrl. \\
\midrule
\texttt{keep\_all} & --- & n/a & --- & n/a \\
\texttt{reset\_m} & +0.021 [+0.015, +0.032] & +0.012 [+0.010, +0.015] & +0.020 [+0.012, +0.027] & +0.008 [-0.005, +0.021] \\
\texttt{reset\_v} & -0.152 [-0.188, -0.087] & -0.020 [-0.060, +0.056] & -0.244 [-0.290, -0.176] & -0.048 [-0.081, +0.013] \\
\texttt{reset\_both} & -0.008 [-0.049, +0.045] & -0.027 [-0.058, -0.011] & -0.043 [-0.063, -0.008] & -0.030 [-0.050, -0.016] \\
\texttt{reset\_t} & -0.207 [-0.255, -0.183] & +0.000 [+0.000, +0.000] & -0.203 [-0.244, -0.130] & +0.000 [+0.000, +0.000] \\
\midrule
\texttt{reset\_v}, shared-keep$^\dagger$ & +0.063 [-0.163, +0.313] & n/a & -0.241 [-0.287, -0.171] & n/a \\
\texttt{reset\_both}, shared-keep$^\dagger$ & +0.333 [+0.268, +0.410] & n/a & +0.251 [+0.235, +0.273] & n/a \\
\bottomrule
\end{tabular}

\end{table}

Table~\ref{tab:interventions} summarises the intervention branches. \keep{} reproduces the late probe of
Section~\ref{sec:phenomenon} bit for bit, and all branches received identical batches and masks.

\textbf{A fresh optimizer state does not recover the gap.} \texttt{reset\_both} with decoupled counters
is the same algorithm as re-creating \texttt{torch.optim.Adam} (Appendix~\ref{app:checks}). Its
difference to \keep{} is \RTResetBothVsKeep{} and \LPResetBothVsKeep{}, far smaller in magnitude than
the late-minus-early gap. In this setting, therefore, the loss of learning ability is not explained by
stale optimizer state alone.

\textbf{First-moment reset: no support for benefit beyond update scale.} \texttt{reset\_m} is the only
correctly corrected intervention with a positive difference to \keep{} on every seed:
\RTResetMVsKeep{} [\RTResetMVsKeepMin, \RTResetMVsKeepMax] (random-teacher, which meets the effect
criterion) and \LPResetMVsKeep{} [\LPResetMVsKeepMin, \LPResetMVsKeepMax] (label-permutation, which does
not). It also roughly doubles the early update norm (\RTResetMUpdTen{} versus \RTKeepAllUpdTen{}).
Against its update-norm control the difference shrinks to \RTResetMVsCtrl{}
[\RTResetMVsCtrlMin, \RTResetMVsCtrlMax] and \LPResetMVsCtrl{}, below the minimum interesting effect.
We read this as \emph{no support for an additional benefit of at least 0.02} beyond matching the update
scale. The random-teacher residual is positive on all three seeds, so we do not claim that it is zero or
that the two arms are equivalent.

\textbf{Counter-only reset slows learning, as its schedule predicts.} \texttt{reset\_t} lowers AUC by
\RTResetTVsKeep{} and \LPResetTVsKeep{}, and its early update norm is \RTResetTUpdTen{} versus \RTKeepAllUpdTen{} (random-teacher).
Consistent with Section~\ref{sec:analysis}(i), its accuracy curve ties its update-norm control
exactly on every seed. The continuous quantities differ only at the level of \TrajLossDiff{} in loss, so
the trajectories are close but not bitwise identical (Appendix~\ref{app:trace}). For this arm the
control is effectively the same algorithm and does not constitute an independent test.

\textbf{Exactly corrected second-moment reset is unstable.} \texttt{reset\_v} raises the mean
first-10-step update norm from \RTKeepAllUpdTen{} to \RTResetVUpdTen{} (random-teacher) and from
\LPKeepAllUpdTen{} to \LPResetVUpdTen{} (label-permutation), and it lowers AUC by \RTResetVVsKeep{} and
\LPResetVVsKeep{}. This is the behaviour \eqref{eq:resetv} allows. At the tensor level, the exploding tensors have shared-keep/decoupled update ratios of at most
\TraceVSkBigMaxRatio{} (\TraceVSkBigNearOne{} of \TraceVSkBigCount{} within 1\% of 1), instead of the
$\eps=0$ prediction \TraceVSkPred{} (Appendix~\ref{app:trace}). By \eqref{eq:ratio}, this is the signature of $\eps$-dominated
coordinates, which is consistent with near-zero current gradients meeting a non-zero kept momentum. We
did not store per-coordinate traces, so we could not verify this at the coordinate level.

\textbf{The largest gains come from a mis-corrected artifact.} \texttt{reset\_both} with the counter
left unchanged improves AUC by \RTResetBothSharedKeepVsKeep{} and \LPResetBothSharedKeepVsKeep{}. This
exceeds the late-minus-early gap in random-teacher. Its early update norm is \RTResetBothSharedKeepUpdTen{}
versus \RTKeepAllUpdTen{}. By Section~\ref{sec:analysis}(ii), under $\eps=0$ this arm is fresh-state Adam
with a deterministic learning-rate multiplier. We did not run a learning-rate-matched control, so we do
not attribute its empirical gain entirely to the learning rate. The observation does show that, in this
setting, the late network's probe learning is highly sensitive to update scale.

\textbf{Descriptive statistics.} All branches update their weights after the intervention. Relative
parameter movement, CKA of hidden features to the starting point, and the accuracy drop on task 49
(Appendix~\ref{app:descriptive}) vary with update scale. For example, the artifact arm moves the
parameters most and has the lowest CKA. We use these only as descriptions, not as evidence that features
were or were not relearned.
